\subsection{Trace polynomial optimization}
% Consider a QKD protocol in prepare and measure scenario, where Alice prepares one of four quantum states randomly with equal probability and sends it to Bob over an untrusted channel where Eve can eavesdrop. Bob does one of three dichotomic measurements with equal probability.
% We try to maximize the probability that Eve can guess Alice input by eavesdropping on the channel. The quantum states have to obey two constraints. First, they cannot be distinguished with more than $G$ probability. Second, they have to reproduce the correlations that Alice and Bob $p(b|x,y) $ find. $G$ and $p(b|xy$) are parameters of the protocol.
% The optimization variables are the four quantum states $\{\rho_1,\rho_2,\rho_3,\rho_4\}$ and free hermitian operator $\sigma$ acting on Bob and Eve's joint quantum system ,the povm for the first coutcome of the three measurements of Bob $\{B_1,B_2,B_3\}$ acting only on Bob's system and Eve's povm $E$ acting only on Eve's system.This can be cast as a trace polynomial optimization problem as follows.

Consider the problem, studied in \cite{Tavakoli2022informationally}, of characterizing correlations in a prepare-and-measure scenario where the information capacity of the ensemble of prepared quantum states is constrained.
The trace polynomial optimization problem reads as follows.
With PCPOP, we reproduce Figure~3 in \cite{Tavakoli2022informationally}.
% \begin{equation}
% \begin{split}
%     p^* &= \underset{\rho_1,\rho_2 , \mathcal{H}}{inf} \: 0.5*tr(\rho_1*E+\rho_2*(1-E)), \\
%     s.t \: & \rho_i - \rho_i^2 ≥ 0, ∀ i ∈ \{1,2,3,4\}\\
%     & σ - 0.25*\rho_i\geq 0, ∀ i ∈ \{1,2,3,4\} \\
%     & -tr(\sigma) ≥ -G \\
%     & tr(\rho_i)=1, ∀ i ∈ \{1,2,3,4\} \\
%     & tr(\rho_i*B_j)=p(1|i,j), ∀ i ∈ \{1,2,3,4\}, j ∈ \{1,2,3\} \\
%     & B_j^2 - B_j = 0, ∀ j ∈ \{1,2,3\} \\
%     & E^2 - E = 0. \\
%     & [B_i,E]=0. ∀ i ∈ \{1,2,3\} \\
% \end{split}
% \end{equation}
% \begin{code}
%     @pcmonoid M ρ[4,0] σ B[3,0] E
%     Projector.([B;E])
%     @comms B E
%     build(M)
%     obj=0.5*(ρ[1]*E + ρ[2]*(1-E))
%     # To impose that the states are mixed states.
%     ge=[i-i*i for i in ρ] 
%     # Following two lines, impose information constraint.
%     ge=vcat(ge,[σ-0.25*ρ[i] for i in 1:4]) 
%     ge_constr=[[-σ,-G]]
%     # To impose normalization of states.
%     eq_constr=[[tr(ρ[i]),1] for i in 1:4]
%     # To impose correlation constraints. where p(1|x,y) is the probability of getting outcome 1 when ALice sent ρ_x and Bob did measurement B_y.
%     eq_constr=vcat(eq_constr,[[tr(ρ[i]*B[j]),p(1|x,y)] for x in 1:4,y in 1:3]) 
%     level="1+B*E+BE*E+BE*B"
%     ov,_,_=npa(obj,level;op_ge=ge,tr_eq=eq_constr,tr_ge=ge_constr,cyclic=true,min=false,normalize=false)    
%     # Notice here we use cyclic=true to do tracial reductions.
% \end{code}

\begin{code}
    G=0.8  #For other values of G, change G and rerun the code. 
    @pcmonoid M ρ[3,0] σ A[2,0]
    Projector.(A)
    build(M)
    PA=[A[1] A[2];1-A[1] 1-A[2]]
    # The following two lines impose information constraint.
    ge = [σ-1/3*ρ[i] for i in 1:3]
    ge_constr = [[-σ,-G]]
    # To impose the constraint that the states are mixed states.
    ge=vcat(ge,[i-i*i for i in ρ])
    # Normalization of states.
    eq_constr = [[ρ[x],1] for x in 1:3]
    wit = -(2*PA[1,1]-1)*ρ[1]-(2*PA[1,2]-1)*ρ[1]-(2*PA[1,1]-1)*ρ[2]+(2*PA[1,2]-1)*ρ[2]+(2*PA[1,1]-1)*ρ[3]
    level = 2 # this is the level of the localising matrices and it will be enough to retrieve the values of the plot (they use level 3 of the principal moment matrix)
    ov,_,_=npa(wit,level;op_ge=ge,tr_eq=eq_constr,tr_ge=ge_constr,cyclic=true,min=false,normalize=false)
    # Notice here we use cyclic=true to do tracial reductions.
\end{code}
\AM{Moi: Maybe you can extend with more examples that solves general state polynomial optimization problems.}